\section{Methods}\label{sec:methods}

\subsection{Simulations}

Simulated tumours were generated with the inferCNV benchmarking toolkit (\texttt{simCNV.R}, https://github.com/broadinstitute/inferCNV\_benchmarking; a copy with explicit seed and output arguments is in the analysis repository) using its mean--variance method, which simulates counts with the same model as the h-spike (rounded, zero-truncated normal counts with dropout; Supplementary Methods) from mean--variance and mean--dropout trends fitted to a set of normal cells. Smart-seq2 parameters were learned from 496 non-malignant cells of synovial sarcoma \citep{jerbyarnon2021} (8,252 genes) and 10x parameters from 262 normal cells (7,210 genes), most likely normal colon epithelial cells from the colorectal cancer atlas of \citet{pelka2021}: of the 111 cells from sample batches present in its public metadata, 80 match by sample and cell barcode, all annotated as epithelial; genes were placed with GENCODE v19 coordinates. Each data set contained 100 simulated normal cells, used as the reference, and tumour cells carrying one to four CNA specifications. Specification 1: chr1 1.5× and chr4 0.5× (whole chromosomes), chr2:10--50 Mb 2× and chr8:10--30 Mb 0.01×. Specification 2 adds chr10 2×; specification 3 adds to specification 2 chr17 0.01× and chr19 0.5×; specification 4 adds to specification 2 chr12 1.5× and chr13 2×. The scenarios were sim1, specification 1 (100 cells); sim2, specifications 1--2 (60 and 40 cells); sim3, specifications 1--3 (100, 60 and 40 cells); and sim4, specifications 1--4 (100, 60, 40 and 30 cells). Each scenario was simulated and analysed with seeds 1--5.

inferCNV was run with \texttt{cutoff\ =\ 1} (Smart-seq2) or 0.1 (10x), \texttt{cluster\_by\_groups\ =\ TRUE}, \texttt{analysis\_mode\ =\ "subclusters"} (Leiden), \texttt{HMM\ =\ TRUE} with \texttt{HMM\_type\ =\ "i6"} or \texttt{"i3"}, \texttt{BayesMaxPNormal\ =\ 0.5}, \texttt{denoise\ =\ TRUE} and 4 threads, with other parameters at their defaults. The results reported are from the master v1.23.0 code, with R's random-number generator set to L'Ecuyer-CMRG and seeded before \texttt{run()}; runs with the release build using \texttt{seed} agreed within 0.006 on every metric.\footnote{The two builds differ in the i3 transition matrix: in master v1.23.0 its diagonal was $1-5t$, so that each row summed to $1-3t$, whereas the release uses $1-2t$. The difference is of the order of $t = 10^{-6}$ per gene and is covered by the agreement just stated.} Calls were scored per cell and gene: the expected state of each gene in each cell came from the specification of the cell's simulated group (neutral for normal cells), and the predicted state from the HMM region covering the gene in the cell's subcluster. Simulated fold changes 0.01, 0.5, 1, 1.5, 2 and 3 map to i6 states 1--6 (0, 1, 2, 3, 4 and $\geq$5 copies); for i3, losses, neutral and gains map to states 1--3. Sensitivity is the fraction of altered positions in tumour cells called in the correct direction; precision the fraction of non-neutral calls in the correct direction; exact-state accuracy the fraction of altered positions assigned the exact i6 state; and the reference call rate the fraction of non-neutral calls in reference cells. Purity is the fraction of tumour cells in subclusters whose majority belongs to their own simulated group, and the adjusted Rand index (ARI; Hubert--Arabie form, computed by the scoring script) measures the agreement between the subcluster partition of the tumour cells and their simulated groups, with 1 for identical partitions and about 0 for a random one.

\subsection{Evaluation against DNA-based copy number}

Data were obtained as in \citet{schmid2025}. MM199 counts were taken from GEO GSE141299 (GSM4200471) \citep{mahdipourshirayeh2022scicnv}. Because the reference annotation used by Schmid et al.~was not deposited, we repeated their cell filtering (more than 3,000 counts, less than 5\% mitochondrial reads) and Louvain clustering (resolution 0.5, 10 principal components) in Seurat, using log normalisation instead of scran. Clusters of plasma cells expressing IgG-λ markers were labelled malignant (1,916 cells), and clusters of B cells and NK/T cells formed the reference (245 cells). BCC06 and BCC06post counts and annotations were taken from GEO GSE123813 \citep{yost2019}; as in Schmid et al., endothelial cells, melanocytes and myofibroblasts of patient su006 formed the reference (33 cells pre-treatment, 280 post-treatment) and the two tumour clusters were the observations (1,403 and 320 cells). Gene positions were from GENCODE v47 (GRCh38), made with \texttt{gtf\_to\_position\_file.py} and restricted to chr1--22, X and Y. Whole-exome copy number segments (GATK) were taken from the Schmid et al.~Zenodo deposit \citep{schmid2025zenodoTruth}.

inferCNV was run with \texttt{cutoff\ =\ 0.1}, \texttt{cluster\_by\_groups\ =\ TRUE}, \texttt{denoise\ =\ TRUE}, \texttt{HMM\ =\ TRUE}, \texttt{HMM\_type\ =\ "i6"}, \texttt{analysis\_mode\ =\ "subclusters"}, \texttt{BayesMaxPNormal\ =\ 0.5}, \texttt{write\_expr\_matrix\ =\ TRUE}, \texttt{seed\ =\ 1} and 4 threads, as in Schmid et al.~except for the seed and thread count. After filtering, 5,740 (MM199), 6,963 (BCC06) and 7,401 (BCC06post) genes remained.

The evaluation re-implements that of Schmid et al.~Truth segments were tiled into 100-kb bins, each bin was assigned the mean of its segment calls (loss = 1, neutral = 2, gain = 3), and bins were called losses below 1.5 and gains above 2.5. The CNV score of a gene was the mean over malignant cells of its filtered i6 call expressed as loss, neutral or gain (representative values 0 and 0.5 mapped to 0, 1 to 1, and 1.5, 2 and 3 to 2), plus 1; the expression score was the per-gene mean of the final residual expression plus 1. Bin scores were the mean over overlapping genes, and only bins covered by inferCNV genes were scored (Schmid et al.~used bins covered by all callers). Metrics were Pearson correlation with the truth, AUC for gains and for losses, partial AUC (sensitivity range from 0 to the sensitivity at score 2) and the maximal unweighted three-class F1 over a grid of loss and gain thresholds, with sensitivity and precision at the optimum. Our re-implementation of the partial AUC reproduced the BCC06 expression-score value reported by Schmid et al.~(0.523). Values for inferCNV v1.10 were taken from the source data of \citet{schmid2025}.

\subsection{Run time}

Run times for Fig.~\ref{fig:4} were measured on a compute cluster with 4 threads per job (wall time and peak resident memory as reported by the scheduler). TNBC1 (10x, 1,097 cells: 302 normal cells as reference and two tumour clones) \citep{gao2021copykat, huang2024xclone, xclonedata} was run in expression-only mode with the Fig.~\ref{fig:4} settings on an 8-core host with load averages of about 40, so its times are upper bounds; step times were taken from the run log. In an earlier run, Leiden subclustering appeared to take 1.6 h; this was caused by calling \texttt{run()} through \texttt{do.call()} with the object inlined in the call, which Seurat's \texttt{ScaleData} then deparses. Calling \texttt{run()} directly avoids this. For the scaling experiment (Supplementary Fig.~\ref{fig:Sscale}), nested down-samples of 500, 1,000, 2,500, 5,000, 10,000 and 20,000 cells were drawn from a 10x breast cancer atlas \citep{wu2021breast} (GEO GSE176078; author cell-type labels): 80\% malignant epithelial cells from the eight tumours with the most such cells, annotated by tumour, and 20\% endothelial cells and cancer-associated fibroblasts as two reference groups, allocated proportionally, with each subset containing the smaller ones. About 8,000 genes passed \texttt{cutoff = 0.1} at every size. Each size was run with i6 and the Bayesian filter (default), i6 without it (\texttt{BayesMaxPNormal = 0}) and i3 without it, with 4 threads on lightly loaded 128-core cluster nodes. Wall time and peak resident memory were taken from \texttt{/usr/bin/time -v}, and step times from the run log.

\subsection{Reference designs}

The PBMC data (10x Genomics PBMC 3k, Cell Ranger 1.1.0, hg19; GENCODE v19 gene positions) were labelled with the Seurat PBMC tutorial workflow (quality filter $200 <$ detected genes $< 2{,}500$ and mitochondrial fraction $< 5\%$; 2,000 variable genes; 10 principal components; clustering at resolution 0.5) and canonical markers, giving 2,625 cells in seven cell types after removal of platelets. Each cell type was split into random halves (\texttt{set.seed(1)}). The observation half (1,314 cells) was the same in every design. The reference was the other half, one group per cell type, or the same cells pooled into one group; no reference (both halves as observations); the reference halves of the lymphoid types (CD4 T, CD8 T, NK and B cells), of the CD4 T cells or of the B cells; 5, 10, 25 or 50 randomly chosen reference cells per type; Kinnex PBMCs from another donor (see \emph{Long-read and spatial data}; up to 500 cells per type, one group per type, genes shared by both data sets); or 200 randomly chosen GTEx v10 whole-blood RNA-seq samples (gene read counts summed by gene symbol, genes shared) \citep{gtex2020}. All PBMC runs used the release defaults described here (\texttt{cutoff\ =\ "auto"}, Leiden subclusters, i6 HMM, \texttt{denoise\ =\ TRUE}, \texttt{BayesMaxPNormal\ =\ 0.5}, \texttt{seed\ =\ 1}); the automatic cutoff was 0.01 (8,671 genes) for the single-cell references, 0.03 with the Kinnex reference and 1 with the GTEx reference. The false-call rate is the fraction of genes $\times$ observation cells in a non-neutral i6 state.

The expected residual of a cell type was computed from the input counts. Counts were scaled to the median library size of the observations, $\log(x+1)$-transformed and averaged over the cells of the type and of each reference group. Over the genes retained by the run, the difference from the reference was zero where the cell type's mean lay within the range of the reference group means, and otherwise its distance to the nearer bound (with one reference group, or none, the difference from that group's mean or from the mean of all cells). It was smoothed along each chromosome with a 101-gene running mean and centred on its median. Its spread is its standard deviation over the genome.

For the automatic reference, all 2,625 cells were given as one annotation group to an experimental two-pass function of the development branch (\texttt{infercnv\_auto\_reference()}), which is not part of the supported interface. In the first pass, inferCNV ran without a reference. A provisional reference was then chosen among the first-pass subclusters, either as those expressing marker panels of immune, endothelial or fibroblast cells at least fourfold more than the other subclusters (as in SCEVAN, \citealp{defalco2023}) or as the expression clusters with the flattest self-centred, smoothed profiles (the principle of CopyKAT, \citealp{gao2021copykat}). In the second pass, with that reference, each observation subcluster was called malignant when both its CNA score (mean squared residual) and its correlation with the average profile of the top 5\% of cells by score exceeded the 99th percentiles of random groups of reference cells of the same size, normal when both were below, and unresolved otherwise \citep{tirosh2016oligo, puram2017}.

The spatial normal-tissue designs used the marker-labelled 32-, 16- and 8-\textmu m bins described in \emph{Long-read and spatial data}, with every class (epithelium, fibroblast, smooth muscle, immune, endothelial) split into random halves, one half of each as a reference group. They were run with the spatial settings and compared with the same bins analysed against the non-epithelial reference groups only.

\subsection{Smart-seq2 tumours with clone and fusion labels}

For the H3K27M gliomas BCH869 and BCH1126 \citep{filbin2018} we used the published RSEM TPM matrix (GEO GSE102130; Single Cell Portal SCP147), and for the synovial sarcomas SS1, SS13 and SS16 \citep{jerbyarnon2021} the published Smart-seq2 counts (GEO GSE131309) with the cell annotations and per-cell SS18--SSX fusion calls of the original study (Single Cell Portal SCP167). Both matrices use hg19 gene symbols, which were placed with the GENCODE v19 gene-order file of inferCNV\_examples (93\% of genes). inferCNV was run with \texttt{cutoff\ =\ 1}, \texttt{cluster\_by\_groups\ =\ TRUE}, \texttt{denoise\ =\ TRUE}, Leiden subclusters, the i6 HMM with the Bayesian filter (\texttt{BayesMaxPNormal\ =\ 0.5}) and \texttt{seed\ =\ 1}. For the synovial sarcoma comparison, each cell's CNV score was the mean squared deviation of its final residuals from 1, its correlation was with the mean residual profile of the 10\% highest-scoring observed cells of the patient, and a cell was classified malignant if either value exceeded the 99th percentile of the patient's reference cells. For BCH869, published clone labels and per-clone copy number were taken from \citet{huang2024xclone}; cell names were matched between data sets (all 492 cells), truth was assigned to genes by their hg19 coordinates, and scores were compared at gene and chromosome-arm level as in the DNA-truth evaluation above.

As a cross-check, we realigned the reads of these cells to GRCh38 (GENCODE v44, CTAT genome library) and processed each cell with STAR-Fusion 1.15.1 \citep{haas2019starfusion}, which runs STAR \citep{dobin2013star} once for both the genome alignment and fusion detection, with gene counts from featureCounts (Subread 2.1.1). SS18--SSX, called when SS18 was fused to any SSX paralog, was detected in 350 synovial sarcoma cells, of which 340 were fusion-positive in the original study (with the same SSX partner in all but three); the original calls included 139 further cells. inferCNV runs on the realigned counts gave similar results: 344 of 349 observed cells with STAR-Fusion-detected SS18--SSX were classified as malignant, and the BCH869 arm-level AUCs were 0.98 for gains and 1.00 for losses.

\subsection{Long-read and spatial data}

All runs used the settings of Fig.~\ref{fig:4} (\texttt{cutoff\ =\ 0.1}, Leiden subclusters, i6 HMM, \texttt{denoise\ =\ TRUE}, 4 threads), with the Bayesian filter (\texttt{BayesMaxPNormal\ =\ 0.5}) for single-cell data and without it for spatial data, and with GENCODE v47 gene positions. Gene-level counts were summed by gene symbol.

\emph{PBMC.} Kinnex 10x 3$'$ PBMC data are from the PacBio public data sets (DATA-Revio-Kinnex-PBMC-10x3p). Genes were quantified with LRAA v0.45.0 in its reference-guided single-cell mode \citep{yu2026lraa}, and gene components without a gene symbol (unannotated loci) were removed. Cell types were those of \citet{yu2026lraa}, with Seurat clusters grouped into CD4 T, CD8 T, CD8 T/NK, NK, B, CD14 and CD16 monocytes and dendritic cells. Reference designs and the false-call rate were as in Fig.~\ref{fig:6}a. The relative depth of a subcluster was its median count total divided by that of its cell type.

\emph{LongBench.} Gene count matrices of the single-cell data sets and the cell-line assignments (merged across platforms) were taken from the LongBench repository \citep{you2025longbench}. Gene identifiers (GENCODE v44) were mapped to symbols with the Cell Ranger feature file of the same data. Each of Kinnex, ONT and Illumina was run without a reference, with all 3,850 assigned cells grouped by cell line. Kinnex was also run with the 9,936 Kinnex PBMCs above as the reference, one group per cell type, on the 18,439 gene symbols shared by the two data sets. CCLE copy number segments (study \texttt{ccle\_broad\_2019}, hg19) were obtained from cBioPortal \citep{cerami2012cbioportal, ghandi2019ccle}. Chromosome arms were defined by centromere positions (hg19 for segments, GRCh38 for genes); the p arms of acrocentric chromosomes and the sex chromosomes were excluded, as were arms with fewer than 20 genes. For each line and arm, the inferCNV value was the mean over genes and cells of the HMM state after Bayesian filtering, minus 3, and the CCLE value was the length-weighted mean segment log$_2$ ratio. For runs without a reference, CCLE values were centred per arm on the mean over the eight lines.

\emph{ccRCC organoids.} Gene count matrices and cell labels (ccRCC or non-ccRCC, assigned by the original authors from expression) of one normal kidney organoid and four ccRCC organoids were taken from Zenodo record 14697181 \citep{zajac2025ccrcc} (2,378 labelled cells; 15,056 gene symbols shared by all samples). All cells of the normal organoid formed the reference, and observations were grouped by organoid and label. A cell was called as having lost 3p if the mean HMM state over its 3p genes after filtering was below 2.5, and as having gained 5q if the mean over 5q was above 3.5; for comparison, a cell's mean 3p residual was called low if it was below the 1st percentile of the normal organoid's cells.

\emph{Visium HD.} Gene count matrices of the Kinnex Visium HD data sets (human breast infiltrating ductal carcinoma, colon adenocarcinoma and tonsil; 2-\textmu m bins) are from the PacBio public data sets \citep{pacbio2026visiumhd}. According to the data set description, the sections (10 \textmu m) are from fresh-frozen tissue purchased from BioIVT: an invasive colon adenocarcinoma (female, 69 years), an infiltrating ductal carcinoma of the breast (female, 70 years) and a chronic hyperplastic tonsil (male, 25 years). Bins were summed into 32-\textmu m spots (16 $\times$ 16 bins), and spots with at least 1,000 counts were kept: in a first analysis with a lower threshold of 300 counts, arm-level calls in spots with fewer than about 1,000 counts tracked depth rather than tissue structure (in the tonsil and in a breast region). Tonsil spots were subsampled to 15,000 before filtering. Spots were clustered with Seurat (log normalisation, 2,000 variable genes, 30 principal components, Louvain resolution 0.5), and each cluster was labelled with the marker set (epithelial, immune, fibroblast, endothelial or smooth muscle) of highest mean module score. In the tumours, epithelial clusters were the observations and fibroblast clusters the reference (colon: one cluster; breast: three, or one in a sensitivity analysis); in the tonsil, a random half of each cluster was the reference for the other half. The CNA score of a spot was the mean squared deviation of its final residuals from 1. Maps show the mean HMM state of each spot's subcluster over the genes of an arm. inferCNV was run with \texttt{cutoff\ =\ "auto"}, which chose 0.03 for each section; arm percentages count spots whose subcluster has a mean state of at least 0.5 above (gain) or below (loss) neutral over the arm.

\emph{Visium HD at several bin sizes.} 10x Visium HD data (formalin-fixed, probe-based) of three colorectal tumours (P1, P2 and P5) and normal adjacent colon (P3) \citep{oliveira2025visiumhd} were downloaded from 10x Genomics as Space Ranger 3.0.0 outputs at 8- and 16-\textmu m bins. 32-\textmu m bins were sums of 2 $\times$ 2 16-\textmu m bins, and the Space Ranger 4.0.1 cell segmentation of P2 provided the segmented cells. In the tumours, 8-\textmu m bins took the labels of the original authors' RCTD deconvolution \citep{cable2022rctd} (singlet bins only), grouped into tumour, normal epithelium, fibroblast, endothelial, smooth muscle and immune classes. A 16- or 32-\textmu m bin took the most common class of its labelled 8-\textmu m bins, and a segmented cell the label of the 8-\textmu m bin containing its centroid. For P3 and the Kinnex sections (8-\textmu m bins as sums of 4 $\times$ 4 2-\textmu m bins), an 8-\textmu m bin was labelled with the marker panel (epithelial, fibroblast, endothelial, smooth muscle or immune) that held the largest share of its counts, if that share was at least twice that of the next panel; larger bins were labelled as above. Observations were tumour bins (10x tumours) or epithelial bins (other sections) in which at least half of the labelled 8-\textmu m bins had that class. The reference comprised bins of the other classes in which at least 75\% did, one group per class. Up to 10,000 observation bins and 1,000 bins per reference class were sampled, and every bin needed at least 100 counts. inferCNV was run in subclusters mode with the i6 HMM, denoising and \texttt{cutoff\ =\ "auto"}, without the Bayesian filter, and with the gene density and null checks set to warn so that every run completed. For each run, the profile was the mean final residual per gene over the observation bins, and the called fraction per gene was the fraction of observation bins in a gain or loss state. Runs of the same section were compared over 5-Mb windows with at least three genes in both runs. Clones were obtained with \texttt{infercnv\_clones()}, and clone assignments were compared on a common 32-\textmu m grid (majority clone per grid square) with the adjusted Rand index. For comparison, all runs were repeated with a cutoff of 0.1 using the earlier code without merging of small subclusters. The null false-call rate of these runs was computed with the procedure of the null check, on their state before the HMM.

\subsection{Software}

inferCNV was run in R 4.4.0 from the Docker image trinityctat/infercnv:1.20.0 (converted to Apptainer), which provides rjags 4-15 with JAGS 4.3.0, HiddenMarkov 1.8-13, Seurat 5.1.0, igraph 2.0.3, caTools 1.18.2, edgeR 4.3.4 and coda 0.19-4.1. The inferCNV package itself was loaded from a separate library placed ahead of the image's (\texttt{R\_LIBS}): master v1.23.0 for the simulations, and the release branch described here, built from master v1.23.0 (package version 1.23.0), for all other analyses. Reads were aligned with STAR 2.7.11a through STAR-Fusion 1.15.1, and counted with featureCounts (Subread 2.1.1). Simulated tumours were generated with \texttt{simCNV.R} from the inferCNV benchmarking toolkit, and all analysis scripts, with the exact settings, are in the analysis repository (Availability).
